% numerical_workflow.tex
%
% Purpose:
%   Give a practical end-to-end workflow for assembling, solving, and validating
%   the current trace-subspace formulation.
%
% Main algorithm:
%   Lists the ordered computational steps from geometry selection through
%   scattering solves or singular-value scans.
%
% Based on:
%   Existing driver scripts and notes/theory/scat_formulation2.md.
%
% Main changes:
%   Converts implementation fragments into a concise checklist suitable for a
%   technical report.
%
% Numerical goal:
%   Provide a reproducible roadmap for future MATLAB validation runs without
%   claiming that this report itself executed them.

\section{Numerical Workflow}

A typical computation proceeds as follows.

\begin{enumerate}
  \item Choose the geometry, transverse quasiperiodicity \(\beta\), and either a
  fixed \(k\) or a scan range for \(k\).

  \item Build the boundary discretization for the center obstacle and, when
  needed, the representative lead obstacles.  Select \(N_{\mathrm{tot}}\) and
  ensure that the quadrature assumptions are compatible with the boundary
  smoothness.

  \item Construct the quasiperiodic Green or proxy/MFS data required by the
  exterior boundary integral operators.

  \item Assemble the center Muller matrix \(\Aqp\) and the center homogeneous
  trace matrices \(H_\Sigma^a,H_\Sigma^b\).

  \item Build the Rayleigh channel data:
  \[
    \beta_m,\qquad \gamma_m,\qquad E=\diag(\ee^{\ii\gamma_mL_0}).
  \]

  \item Assemble each lead-cell scattering matrix \(S_{\mathrm{cell}}^\pm\).

  \item Solve the Bloch generalized eigenproblem
  \[
    \Asc z=\lambda\Bsc z
  \]
  for each lead.

  \item Convert Bloch eigenvectors to wall Cauchy traces and select incoming and
  outgoing half-lead trace subspaces.  Apply the negative-port Neumann sign
  convention when constructing \(N_{\mathrm{out}}^-\).

  \item Build the density far-field extractors \(F_-\) and \(F_+\).

  \item Assemble the global trace-matching matrix \(\Atr\).

  \item For forced scattering, choose valid incoming coefficients
  \(p_{\mathrm{in}}^\pm\) and solve
  \[
    \Atr x=b_{\mathrm{in}} .
  \]

  \item For an eigenvalue or TEP-type scan, compute
  \[
    \sigmin(\Atr(k,\beta))
  \]
  across the parameter range and refine candidate dips.

  \item Validate convergence in \(N_{\mathrm{tot}}\), \(M\), proxy parameters,
  and mode-selection thresholds.  Check BIE residuals, port residuals,
  outgoing-subspace stability, and field plots.
\end{enumerate}

This workflow is intentionally modular.  It allows the same lead-cell
calculation to support either a forced scattering solve or a homogeneous
singular-value scan, and it separates errors caused by boundary quadrature,
quasiperiodic periodization, Rayleigh truncation, and Bloch-mode selection.
